\section{The matched plane and its traces}
\label{sec:characteristic-plane}

The trace construction applies to the matched process of
Proposition~\ref{prop:isotropic-reduction} for any least aspect $m>0$.
In the proof of Theorem~\ref{thm:classical-growth}, this is the
alternative remaining under $\beta>\gamma$ after
Proposition~\ref{prop:isotropic-exclusion}.  Its saturated apertures have shapes
\[
 (b(u),a(u))=(b(u),b(u)+mu),\qquad 0<u<1,
\]
with horizon $L=1+m$, a Lipschitz function $b$, and
$0\le-b'\le1+m$ almost everywhere.  Their labels use one common time
filtration and the prefixes of one orientation, as supplied by
Proposition~\ref{prop:isotropic-reduction}.

\subsection{The information profile and the contact relation}

Work in one bounded slope chart, using the row frame from the end of
Proposition~\ref{prop:isotropic-reduction}.  The particle $T$ carries
$(V,A)$.  At label time $v$ put $E_v=(C_v,R_{mv})$ and set
\begin{equation}\label{eq:trace-variables}
 X=V_1,\qquad Y=A_1+C_vX,\qquad
 W=V_2-R_{mv}X,\qquad
 Z=A_2+C_vV_2-R_{mv}Y.
\end{equation}
For data depths $\mathbf L=(L_X,L_Y,L_W,L_Z)$ let $O(\mathbf L,v)$ record $E_v$ and
these four variables at their respective grid precisions.  We allow
the depths to vary on compact subsets of
\begin{equation}\label{eq:trace-domain}
 \begin{aligned}
 L_Y-L_X&\le v,& L_W-L_X&\le mv,\\
 L_Z-L_W&\le v,& L_Z-L_Y&\le mv,
 \qquad 0<v<1.
 \end{aligned}
\end{equation}
They may range over any fixed bounded extension of the depths of the
apertures; a negative depth gives only boundedly many bins.

When $m=1$, impose $L_Y=L_W=L_U$ at this point and regard
$U=(Y,W)$ as one datum of capacity two, with
$\mathbf L=(L_X,L_U,L_Z)$.  The profile and all the trace
statements below are then defined on this grouped domain, and the middle
datum has one joint trace rate.

The finite matched apertures occupy a particular plane in the depth
domain.  At time $u$, an aperture with common depth $i$ has velocity
depths $(i,i+mu)$ and position depths $(i+u,i+(1+m)u)$.  Define
\begin{equation}\label{eq:matched-plane}
 \begin{aligned}
 \Pi_m(i,u)&=(i,i+u,i+mu,i+(1+m)u,u), &&m\ne1,\\
 \Pi_1(i,u)&=(i,i+u,i+2u,u), &&m=1.
 \end{aligned}
\end{equation}
The coordinate orders are $(L_X,L_Y,L_W,L_Z,v)$ and
$(L_X,L_U,L_Z,v)$, respectively.  The finite matched observation is
$O(\mathbf L,u)$ with $(\mathbf L,u)=\Pi_m(i,u)$.
Every inequality in \eqref{eq:trace-domain} is an equality on this
plane.  The part corresponding to admissible apertures is
\begin{equation}\label{eq:matched-aperture-region}
 \mathcal A_m=\{(i,u):0<u<1,\quad i\ge0,\quad i+(1+m)u\le L\}.
\end{equation}
The path observations are the points $i=b(u)$, up to the zero-cost
grid choices already established for the common row frame.
The bounded extension of the data depths gives a neighborhood within
this plane across the boundary of $\mathcal A_m$.  The full depth domain
still has the one-sided faces in \eqref{eq:trace-domain}.

We now make one homogeneous choice before defining the limiting
functions.  Start with a countable dense family of admissible depth
configurations and the countable family of all finite tuples formed
from their observations, including the versions conditional on $T$.
Also include the score observations at dyadic path times.  At each
stage use only a finite initial list and finitely many path times.
Apply the cap regularization of
Lemma~\ref{lem:classical-cap-regularization}, then the uniform bins of
Lemma~\ref{lem:homogenization}.  By
Lemma~\ref{lem:classical-splitting}, select typical components whose
finitely many score deficits lie under one deterministic vanishing
envelope.  Let the lists grow and the bin
meshes shrink slowly with the precision.  The diagonal of
Lemma~\ref{lem:profile-diagonal}, as in
Remark~\ref{rem:classical-profile-inheritance}, gives one sequence of
finite laws $p_\eps$ carrying all the listed homogeneous profiles.
Every listed path observation is saturated along this sequence in the limit.
Only after this sequence is chosen do we define
\[
 f_\eps(\mathbf L,v)=\Mut^{p_\eps}(T;O(\mathbf L,v)),\qquad
 S_\eps(v)=\Ent^{p_\eps}(C_v\mid T).
\]
Lemma~\ref{lem:profile-diagonal} and the box comparisons below give a
limiting profile $f$ on \eqref{eq:trace-domain}; write $S$ for the
limiting time profile.  The finite root suprema are bounded, so pass to
a further subsequence and write $\tau$ for the limit of $\tau_L(T)$.
The limiting score along the path is continuous: $b,S$ are Lipschitz,
and the box comparisons give continuity of $f$ along
$\Pi_m(b(u),u)$.  Saturation at the dense dyadic times therefore
extends to every $0<u<1$.  The
positive constant in \eqref{eq:orientation-growth} is obtained for
this selected component by the least-aspect argument; it need not be
the constant before the split.  The symbols $f,S,\tau$ below refer to
this component, whose profiles may differ from those before the split.
Later, a trace point and a positive block length
will be chosen for this fixed $f$ before the required real-depth bins
and finite precision are selected.

The information of a matched aperture is therefore the restriction
\begin{equation}\label{eq:matched-plane-profile}
 \Phi(i,u)=f\bigl(\Pi_m(i,u)\bigr).
\end{equation}
The weight of the aperture at $(i,u)$ is $2i+(1-\gamma)mu$.
Rearranging the universal score bound and using saturation on the
path gives
\begin{equation}
 \Phi(i,u)\ge
 2i+(1-\gamma)mu+2S(u)-\beta u-\tau,
 \qquad\text{with equality when }i=b(u).
 \label{eq:closure-obstacle}
\end{equation}
We call this equality the contact relation.  The local projection
arguments will constrain the information rates at almost every plane
point with $i>0$, while saturation is known on this particular contact
curve.  We therefore need both local affine approximations on the
plane and control of how the rates cross the curve.

At a differentiability point of $b$ put $p=-b'(u)$.  A direction
$V_c=\partial_u-c\partial_i$ preserves the data depth $i+cu$, and
its crossing factor at the contact curve is $p-c$.  Here the clocks
are $0,1,m,1+m$, with the two middle clocks grouped when $m=1$.
Figure~\ref{fig:characteristic-plane} shows these relations.
\begin{figure}[H]
\centering
\begin{tikzpicture}[font=\small,>=Stealth,
  direction/.style={->,line width=0.65pt},
  guide/.style={figuregray,dashed,line width=0.5pt}]
  \begin{scope}[x=1.5cm,y=3.6cm]
    \node[font=\small\bfseries] at (1.50,1.22) {Admissible resolutions};
    \fill[figuregray!9] (0,0) -- (3,0) -- (0,1) -- cycle;
    \draw[figuregray,line width=0.65pt] (3,0) -- (0,1)
      node[pos=0.51,above=3pt,sloped,font=\footnotesize] {$i+3u=3$};
    \draw[->] (0,0) -- (3.30,0) node[right] {$i$};
    \draw[->] (0,0) -- (0,1.13) node[above] {$u$};
    \node[below left] at (0,0) {$0$};
    \node[below] at (3,0) {$3$};
    \node[left] at (0,1) {$1$};
    \draw[figureblue,line width=0.95pt] (1.5,0) -- (0,1);
    \node[text=figureblue,anchor=west,font=\footnotesize]
      at (1.42,0.18) {$i=b(u)$};
    \draw[figureblue,line width=0.4pt] (1.40,0.18) -- (1.305,0.13);
    \node[font=\normalsize] at (1.34,0.46) {$+$};
    \node[font=\normalsize] at (0.31,0.46) {$-$};
    \fill (0.825,0.45) circle[radius=1.4pt];
    \draw[guide] (0.575,0.34) rectangle (1.075,0.56);
    \node[anchor=north,font=\footnotesize] at (1.5,-0.16)
      {$m=2,\quad L=3,\quad b(u)=\tfrac32(1-u)$};
  \end{scope}
  \begin{scope}[shift={(10.0,1.30)}]
    \node[font=\small\bfseries] at (0,3.09) {Directions at a contact point};
    \draw[->,figuregray!70] (-2.65,0) -- (2.08,0)
      node[right,text=black] {$\Delta i$};
    \draw[->,figuregray!70] (0,-1.55) -- (0,2.45)
      node[above,text=black] {$\Delta u$};
    \draw[figureblue,dashed,line width=0.8pt]
      (-2.46,1.64) -- (1.98,-1.32);
    \node[text=figureblue,anchor=north east,font=\footnotesize]
      at (1.83,-1.33) {tangent: $\Delta i=-p\,\Delta u$};
    \draw[direction] (0,0) -- (0,2.17);
    \node[anchor=west,font=\footnotesize] at (0.12,1.91) {$c=0$};
    \draw[direction] (0,0) -- (-1.54,1.54);
    \node[anchor=south east,font=\footnotesize] at (-1.58,1.64) {$c=1$};
    \draw[direction] (0,0) -- (-1.95,0.975);
    \node[anchor=south east,font=\footnotesize] at (-2.04,1.04) {$c=m=2$};
    \draw[direction] (0,0) -- (-2.09,0.6967);
    \node[anchor=north east,font=\footnotesize] at (-2.18,0.64) {$c=1+m=3$};
    \fill (0,0) circle[radius=1.4pt];
    \node[font=\normalsize] at (0.80,0.75) {$+$};
    \node[font=\normalsize] at (-0.85,-0.73) {$-$};
    \node[font=\footnotesize,anchor=west] at (0.54,1.39)
      {$V_c=\partial_u-c\partial_i$};
  \end{scope}
\end{tikzpicture}
\caption{The resolution triangle $i\ge0$, $u\ge0$,
$i+(1+m)u\le L$ and the characteristic directions $di/du=-c$,
where $c\in\{0,1,m,1+m\}$.
Each direction $V_c=\partial_u-c\partial_i$ preserves the
corresponding data depth $i+cu$.
The explicit example on the left has $m=2$ and contact slope
$p=-b'(u)=3/2$; its two axes use different drawing scales.
The right panel uses equal scales for the local increments.
The signs mark $\pm(i-b(u))>0$.
At a transverse contact, $p\ne c$, the crossing factor is
$V_c(i-b(u))=p-c$.  Existence of one-sided traces and the associated
jump identities use the hypotheses of Lemma~\ref{lem:curve-traces}.}
\label{fig:characteristic-plane}
\end{figure}


Ambient almost-everywhere differentiability does not solve the local
problem: the entire matched plane lies on the boundary of
\eqref{eq:trace-domain}, and could lie in the ambient exceptional set.
We will obtain the needed traces from the signs of the mixed
derivatives.  Their affine approximation will produce finite entropy
blocks for Section~\ref{sec:rate-transfer}; their measure identities
will later pass those rate constraints to the contact curve.

Here are the properties of $f$ supplied by the finite observations.
It is locally Lipschitz and nondecreasing in every coordinate; its data
derivatives lie between zero and the capacity of that coordinate.  Its
label derivative is nonnegative and locally bounded.  Every mixed
derivative involving two different coordinates, including a data
coordinate and $v$, is a nonpositive distribution in the interior.
We verify the signs, especially the one involving the label.

Advance $v$ to $v'>v$.  Write $a=C_{v'}-C_v$ and
$r=R_{mv}-R_{mv'}$.  Nested labels give
$|a|\lesssim\eps^v$, $|r|\lesssim\eps^{mv}$, and the change of data is
\begin{equation}\label{eq:trace-frame-change}
 X'=X,\quad Y'=Y+aX,\quad W'=W+rX,\quad
 Z'=Z+aW+rY+arX.
\end{equation}
In a data box, the uncertainty in $Y'$ is bounded by a constant times
$\eps^{L_Y}+\eps^{v+L_X}$, and that in $W'$ by a constant times
$\eps^{L_W}+\eps^{mv+L_X}$.  The uncertainty in $Z'$ is bounded by
\[
 C\bigl(\eps^{L_Z}+\eps^{v+L_W}
       +\eps^{mv+L_Y}+\eps^{(1+m)v+L_X}\bigr).
\]
The four inequalities in \eqref{eq:trace-domain} bound these by the
respective target widths.  The inverse shear has the same bounds.
Consequently, conditional on $E_{v'}$, the old and new data descriptions
at fixed admissible depths determine each other with bounded ambiguity.
This assertion also holds at each corner of any coordinate rectangle
contained in the domain.

At fixed labels, a data refinement is determined by the particle and
labels.  Its increase in particle information is therefore exactly its
conditional entropy, up to the bounded bin ambiguities.  A further data
refinement can only decrease this conditional entropy.  For a label
refinement, express both ends of the data refinement in the old frame
conditional on $E_{v'}$, using \eqref{eq:trace-frame-change}.  The two
information increments being compared are then
\[
 \Ent(D_{\rm fine}\mid E_v,D_{\rm coarse})
 \quad\hbox{and}\quad
 \Ent(D_{\rm fine}\mid E_{v'},D_{\rm coarse}),
\]
up to $O(1/\log(1/\eps))$; the latter is no larger than the former.
Here $D_{\rm coarse}$ includes every other data coordinate.  Thus each
mixed rectangle difference of $f_\eps$ is nonpositive up to an error
vanishing as $\eps\downarrow0$.  Passing to the limiting profile proves
the mixed signs, first for rectangle differences and then for
distributions, by integrating the rectangle inequalities against
nonnegative smooth test functions and shrinking the rectangles.

Refining a scalar datum by $d$ costs at most
$d+O(1/\log(1/\eps))$ in normalized entropy; for $U$ the bound is
$2d+O(1/\log(1/\eps))$.  Advancing the labels by $d$ costs at most
$(1+m)d+O(1/\log(1/\eps))$.  The old data are recoverable from the new
observation, and the change of data conditional on the new labels has
bounded ambiguity.  These facts give monotonicity and the asserted
Lipschitz bounds.  For two nearby points on flat faces, move first to a
common slightly finer label so that coordinatewise comparison is
admissible.  This proves the same bounds up to those faces.  The argument
uses vector conditional entropy unchanged when $Y,W$ are grouped.

These comparisons also upgrade the dense diagonal to locally uniform
convergence.  Fix a compact subset $K$ of \eqref{eq:trace-domain} and a
compact relative neighborhood $K^+$ of it.  Uniformly for nearby
$z,z'\in K^+$, the box bounds give
\[
 |f_\eps(z)-f_\eps(z')|
 \le C_K\|z-z'\|_1+O_K(1/\log(1/\eps)).
\]
The same estimate holds at the faces: the common finer label needed
there advances time by $O_K(\|z-z'\|_1)$, and stays in the domain for
nearby points because $K^+$ has label times bounded away from the
endpoints.  For a fixed sufficiently small $\delta>0$, choose a finite
$\delta$-net $z_1,\ldots,z_N$ for $K$ from the dense family in $K^+$.  The same
Lipschitz bound for $f$ then gives
\[
 \sup_{z\in K}|f_\eps(z)-f(z)|
 \le \max_{1\le j\le N}|f_\eps(z_j)-f(z_j)|
       +2C_K\delta+O_K(1/\log(1/\eps)).
\]
First let $\eps\to0$ along the selected sequence, using convergence at
these finitely many net points, and then let $\delta\downarrow0$.
Thus $f_\eps\to f$ locally uniformly on the full closed-sided domain,
including its faces.

\subsection{An abstract trace theorem}

We now isolate the consequence of these mixed signs for a plane whose
data depths move at distinct speeds.  The grouped middle datum at $m=1$
is what makes the speeds distinct in that case.

Let $n\ge2$, let $c_1<\cdots<c_n$ be distinct real numbers, and let
$\mathcal E$ be a nonempty set of ordered pairs $(j,k)$ with $j<k$.
For an open interval $I$ define the closed-sided region
\[
 \mathcal D=\{(\mathbf L,v)\in\RR^n\times I:
       L_k-L_j\le(c_k-c_j)v\text{ for }(j,k)\in\mathcal E\}.
\]
Its interior is obtained by making these inequalities strict.  Define
\begin{equation}\label{eq:trace-plane-cone}
 \begin{split}
 \Pi(i,u)&=(i+c_1u,\ldots,i+c_nu,u),\\
 \mathcal K&=\{\xi\in\RR^{n+1}:
       \xi_k-\xi_j\le(c_k-c_j)\xi_v
                      \text{ for }(j,k)\in\mathcal E\},\\
 V_j&=\partial_u-c_j\partial_i.
 \end{split}
\end{equation}
The cone $\mathcal K$ has nonempty interior and is invariant under
translation by either tangent vector
$e=(1,\ldots,1,0)$ or $t=(c_1,\ldots,c_n,1)$.
In particular, for $h\ge0$, $\Pi(i,u)+h\xi$ belongs to $\mathcal D$ for
$\xi\in\mathcal K$, provided its label time belongs to $I$.

\begin{theorem}[Characteristic-plane traces]\label{thm:characteristic-trace}
Suppose $f$ is locally Lipschitz on $\mathcal D$, and, in its interior,
\[
 0\le f_j\le d_j,\qquad 0\le f_v\le d_v,
 \qquad f_{jk}\le0\ (j\ne k),\qquad f_{jv}\le0
\]
in the distributional sense.  The derivative bounds may instead be local
bounds on each compact neighborhood under consideration; in that
case all conclusions concerning boundedness are local.  Then there
are, up to null sets, uniquely determined locally bounded fields $a_1,\ldots,a_n,D$ on
$\RR\times I$, with $0\le a_j\le d_j$ and $0\le D\le d_v$, having the
following properties.

\begin{enumerate}
\item \emph{Strong traces.} For almost every interior offset $\eta\in\mathcal K^\circ$, the
restrictions of the ambient gradients to $\Pi(\cdot)+\eta$ are defined
almost everywhere.  As $\eta\to0$ through such offsets, they converge
strongly in $L^1_{\rm loc}(di\,du)$ to $(a_1,\ldots,a_n,D)$.
For $Q\Subset Q'\Subset\RR\times I$ this convergence has the bound
\begin{equation}\label{eq:trace-quantitative-slices}
 \|f_\alpha(\Pi(\cdot)+\eta)-A_\alpha\|_{L^1(Q)}
 \le C_Q|\eta|+\omega_{\alpha,Q'}(C_Q|\eta|),
 \qquad A=(a_1,\ldots,a_n,D),
\end{equation}
where $\omega_{\alpha,Q'}(r)\to0$ is the $L^1$ translation modulus of
$A_\alpha$ on $Q'$.  The constant depends only on the local derivative
bounds, the clock speeds, and the two patches.

\item \emph{Measure identities.} There exist locally finite nonnegative measures
$\nu_{jk}=\nu_{kj}$, $j\ne k$, and $\nu_{jv}$ on the plane such that
\begin{equation}\label{eq:trace-mixed-measures}
 V_j a_j=-\nu_{jv}-\sum_{k\ne j}(c_k-c_j)\nu_{jk},
 \qquad \partial_iD=-\sum_j\nu_{jv}.
\end{equation}
For $F=f\circ\Pi$, the weak derivatives are
\begin{equation}\label{eq:trace-plane-derivatives}
 F_i=\sum_j a_j,\qquad F_u=\sum_j c_ja_j+D.
\end{equation}
In particular,
\begin{equation}\label{eq:trace-distribution-inequality}
 V_n a_n\ge\partial_iD
\end{equation}
as measures.  With $q_j=1+c_n-c_j$, the full flux identity is
\begin{equation}\label{eq:trace-measure-flux}
 -\sum_jq_jV_ja_j
 =\sum_{j<k}(c_k-c_j)^2\nu_{jk}+\sum_jq_j\nu_{jv}.
\end{equation}

\item \emph{Affine approximation.} There is a sequence $h_\ell\downarrow0$ and a single null set
$N\subset\RR\times I$ such that, for every $(i,u)\notin N$ and every
finite $R$,
\begin{equation}\label{eq:trace-affine-blowup}
 \lim_{\ell\to\infty}\sup_{\xi\in\mathcal K,\ |\xi|\le R}
 \left|
 \frac{f(\Pi(i,u)+h_\ell\xi)-f(\Pi(i,u))}{h_\ell}
      -\sum_j a_j(i,u)\xi_j-D(i,u)\xi_v
 \right|=0.
\end{equation}
Only sufficiently large $\ell$, for which the displayed label times
belong to $I$, are used.  The sequence may depend on $f$.
\end{enumerate}
\end{theorem}

\begin{proof}
All estimates are local.  Fix nested bounded open plane patches
$Q\Subset Q'\Subset Q''$ and a nonnegative smooth function $\chi$
supported in $Q'$ and equal to one on $Q$.  All small offsets below keep
$\Pi(Q'')+\eta$ in the fixed neighborhood where the derivative bounds
hold.  Smooth $f$ by convolution at a radius smaller than the distance
of the offset patch from the boundary.  The mixed signs and first
 derivative bounds persist.  We first work with this smooth function.

On any parallel interior plane the chain rule gives
\begin{equation}\label{eq:trace-smooth-flux}
 \begin{split}
 \sum_jq_jV_j f_j(\Pi+\eta)
 &=\sum_{j<k}(q_j-q_k)(c_k-c_j)f_{jk}(\Pi+\eta)
       +\sum_jq_jf_{jv}(\Pi+\eta)\\
 &=\sum_{j<k}(c_k-c_j)^2f_{jk}(\Pi+\eta)
       +\sum_jq_jf_{jv}(\Pi+\eta).
 \end{split}
\end{equation}
Let $M$ bound the absolute values of the first derivatives on this
neighborhood.  Multiplication by $-\chi$ and integration by parts gives
\begin{equation}\label{eq:trace-flux-bound}
 \begin{split}
 &\sum_{j<k}(c_k-c_j)^2
    \int\chi(-f_{jk})(\Pi+\eta)
       +\sum_jq_j\int\chi(-f_{jv})(\Pi+\eta)\\
 &\hspace{35mm}\le M\sum_jq_j\|V_j\chi\|_{L^1}.
 \end{split}
\end{equation}
Every integrand on the left is nonnegative.  Since the speeds are
distinct, this controls each mixed derivative separately, uniformly in
the offset and smoothing radius.  For example, the sum of all their
masses on $Q$ is at most
\[
 C_Q=\frac{M\sum_jq_j\|V_j\chi\|_1}
 {\min\{1,\min_{j<k}(c_k-c_j)^2\}}.
\]
Changing $C_Q$ to accommodate larger patches causes no difficulty.

We next construct the strong traces.  Write $\alpha$ for one ambient
coordinate, either a datum or $v$.  Consider interior offsets with
$\eta_\alpha=0$.  Their cone is convex and nonempty: for a data
coordinate $j$ one may take $\eta_v>0$ and all data offsets zero; for
$\alpha=v$ one may take $\eta_k=-c_k$ and $\eta_v=0$.
If $\eta_\alpha=\eta'_\alpha=0$, the fundamental theorem of calculus
along their segment, followed by \eqref{eq:trace-flux-bound}, gives
\begin{equation}\label{eq:trace-own-coordinate}
 \begin{split}
 &\|f_\alpha(\Pi+\eta)-f_\alpha(\Pi+\eta')\|_{L^1(Q)}\\
 &\quad\le\sum_{\beta\ne\alpha}|\eta_\beta-\eta'_\beta|
       \int_0^1\int_Q
       |f_{\alpha\beta}(\Pi+(1-s)\eta+s\eta')|\,di\,du\,ds
 \le C_Q|\eta-\eta'|.
 \end{split}
\end{equation}
Only mixed derivatives occur in this estimate.

For completeness, this slice estimate passes to the nonsmooth function
without choosing an arbitrary restriction of an $L^\infty$ field.  On
compact interior sets the smoothed gradients converge to the weak
gradients in $L^1$.  The linear map
\[
 (i,u,\eta)\longmapsto\Pi(i,u)+\eta,
 \qquad \eta_\alpha=0,
\]
has full rank.  Fubini's theorem and a subsequence therefore give
$L^1(Q)$ convergence of the smoothed restrictions for almost every such
offset.  Exhaust the allowable offset sets and patches countably.  For
any two offsets in the resulting full-measure set, their compact
connecting segment is interior, so \eqref{eq:trace-own-coordinate}
passes to the limit.  It makes the restrictions an $L^1(Q)$ Cauchy family
as their offsets tend to zero.  The limits on different patches agree;
call the resulting field $A_\alpha$.  Letting $\eta'\to0$ gives
\[
 \|f_\alpha(\Pi+\eta)-A_\alpha\|_{L^1(Q)}\le C_Q|\eta|
 \quad\text{when }\eta_\alpha=0.
\]
The derivative bounds pass to this limit.

For a general data offset remove its $j$th component by replacing
$\eta$ with $\eta-\eta_je$ and translating $i$ by $\eta_j$.
For the label derivative replace $\eta$ with $\eta-\eta_vt$ and
translate $u$ by $\eta_v$.  Cone membership is unchanged by these
operations.  Translation continuity in $L^1$ now proves
\eqref{eq:trace-quantitative-slices}, for almost every general offset.
It also proves uniqueness of the traces.  No pointwise modulus of
continuity of the gradients has been used.

Choose interior offsets tending to zero for which all the gradient
restrictions converge strongly on the countable exhaustion of patches.
For each offset choose a still smaller convolution radius so that the
smoothed restrictions have the same limits.  The nonnegative measures
on each slice with densities $-f_{jk}$ and $-f_{jv}$ have locally bounded
masses by \eqref{eq:trace-flux-bound}.  Successive weak measure limits,
with a diagonal choice on the patches, give $\nu_{jk}$ and $\nu_{jv}$.
The slice identities
\[
 V_jf_j=f_{jv}+\sum_{k\ne j}(c_k-c_j)f_{jk},
 \qquad \partial_if_v=\sum_jf_{jv}
\]
pass to the limit against smooth test functions, proving
\eqref{eq:trace-mixed-measures} and \eqref{eq:trace-measure-flux}.
The potentials on these slices converge locally uniformly to $F$,
whereas their tangential derivatives converge strongly in $L^1$.
This gives \eqref{eq:trace-plane-derivatives}.  Finally,
\[
 V_na_n-\partial_iD
 =\sum_{k<n}(c_n-c_k)\nu_{nk}+\sum_{j<n}\nu_{jv}\ge0,
\]
which proves \eqref{eq:trace-distribution-inequality}.
The measures used to represent the mixed derivatives need not be
unique; the trace fields and the distributional conclusions are unique.

It remains to prove the assertion about this particular plane.
For a bounded base patch $Q$ and an integer $R\ge1$, define
\[
 e_{h,Q,R}(i,u)=\int_{\mathcal K\cap B_R}
       |\nabla f(\Pi(i,u)+h\xi)-A(i,u)|\,d\xi.
\]
The boundary of the cone has ambient measure zero and is irrelevant to
this integral.  The strong slice convergence, Fubini's theorem and the
boundedness of the gradients show that
\begin{equation}\label{eq:trace-integrated-gradient}
 \int_Qe_{h,Q,R}\longrightarrow0\qquad(h\downarrow0).
\end{equation}
More quantitatively the integral is bounded by a constant depending on
$R,Q$ times $h+\sum_\alpha\omega_{\alpha,Q'}(C_Rh)$.
Choose $h_\ell\downarrow0$ so that the integrals in
\eqref{eq:trace-integrated-gradient} are at most $2^{-\ell}$ for the
first $\ell$ members of a countable exhaustion of patches and radii.
The summability of these integrals implies
$e_{h_\ell,Q,R}(i,u)\to0$ outside one null set.

Fix a point outside that set.  The rescaled potentials
\[
 g_\ell(\xi)=
 \frac{f(\Pi(i,u)+h_\ell\xi)-f(\Pi(i,u))}{h_\ell}
\]
are uniformly Lipschitz on each bounded part of $\mathcal K$ and vanish
at zero.  Every locally uniformly convergent subsequence has weak
gradient equal to the constant $A(i,u)$ in $\mathcal K^\circ$, by the
mean gradient convergence just proved.  This interior is connected.
Thus the limit is $A(i,u)\cdot\xi$ there, and continuity gives the same
identity on the closed cone, including its vertex and faces.
Compactness of the uniformly Lipschitz functions and uniqueness of
this limit imply convergence of the whole chosen sequence, as stated
in \eqref{eq:trace-affine-blowup}.
\end{proof}

\subsection{Finite approximation after selecting the block scale}

The next consequence specifies exactly what is required of the finite
entropy systems.  Its exceptional set is a set of base points on the
plane; probabilistic trimming of conditional cells is a separate use
of Lemma~\ref{lem:homogenization} and
Theorem~\ref{thm:finite-conditioning}.

\begin{corollary}[Quantified finite profile approximation]
\label{cor:finite-characteristic-trace}
Assume Theorem~\ref{thm:characteristic-trace}, and let $f_\eps\to f$
locally uniformly on $\mathcal D$ along the sequence of finite systems
under consideration.  Fix a bounded measurable base patch $Q$ of
positive measure with compact closure in $\RR\times I$, a radius
$R\ge1$, an error $\eta>0$, an exceptional proportion $0<\tau<1$,
and a maximum block length $h_*>0$.  There exist
\[
 0<h<h_*,\qquad G\subset Q,\qquad |Q\setminus G|\le\tau|Q|,
\]
and a compact neighborhood $K_0\subset\mathcal D$ containing all
$\Pi(q)+h\xi$ with $q\in\overline Q$, $\xi\in\mathcal K\cap B_R$,
such that the following holds.  If
\begin{equation}\label{eq:trace-finite-tolerance}
 \|f_\eps-f\|_{L^\infty(K_0)}\le\eta h/4,
\end{equation}
then, simultaneously for every $q\in G$ and every
$\xi\in\mathcal K\cap B_R$,
\begin{equation}\label{eq:trace-finite-affine}
 |f_\eps(\Pi(q)+h\xi)-f_\eps(\Pi(q))-hA(q)\cdot\xi|
 \le\eta h.
\end{equation}
In particular, for every finite family of such displacements, every
information increment between two of them differs from its affine
value by at most $2\eta h$.  An alternating sum of $r$ profile values
whose affine contributions cancel has absolute value at most
$r\eta h$.

All dependencies have the order
\[
 f,\ Q,R,\eta,\tau,h_*\quad\longrightarrow\quad h,G,K_0
 \quad\longrightarrow\quad\eps_0.
\]
For every member of the converging finite sequence with
$\eps<\eps_0$, \eqref{eq:trace-finite-tolerance} holds.  If finitely many
entropy identities also incur a total normalized ambiguity error at
most $B/\log(1/\eps)$, one may, after choosing $h$, require in addition
\[
 \log(1/\eps)\ge B/(\eta h).
\]
Each such combined identity then has at most one additional $\eta h$
of error.  More generally one may prescribe any sequences
$\eta_\ell,\tau_\ell\downarrow0$ and $h_{*,\ell}\downarrow0$ and
select successively $h_\ell<h_{*,\ell}$, then the inner approximation
thresholds.  If $\sum_\ell\tau_\ell<\infty$, almost every point of $Q$
belongs to all but finitely many of the resulting $G_\ell$.
\end{corollary}

\begin{proof}
On $Q$, the supremum error in \eqref{eq:trace-affine-blowup} tends to
zero almost everywhere along the chosen sequence and is uniformly
bounded by the Lipschitz bounds times $R$.  Choose a term $h<h_*$
small enough that the subset on which this error exceeds $\eta/2$
has measure at most $\tau|Q|$, and that all shifted label times stay
inside $I$.  Its complement is $G$.  A compact neighborhood $K_0$
exists because $Q$ and the displacement set are bounded.  The two
profile approximation errors in the difference in
\eqref{eq:trace-finite-affine} total at most $\eta h/2$, proving that
inequality.  Subtracting increments, or summing the individual errors,
gives the finite-family assertions.  Uniform convergence supplies
$\eps_0$ after $h$ has been fixed.  The ambiguity condition and the
last exceptional-set assertion follow respectively by addition of
errors and summability of the measures of $Q\setminus G_\ell$.
\end{proof}

There is no uniform choice of $h$ for all Lipschitz profiles in this
corollary.  The translation moduli in
\eqref{eq:trace-quantitative-slices}, and hence the block scale, depend
on the limiting profile.  At a fixed nonexceptional plane point one
may equivalently use its sequence in
\eqref{eq:trace-affine-blowup}, choose a block from that sequence, and
then make the finite-system errors small compared with that block.

For the profile \eqref{eq:trace-domain}, the clocks are
$c_X=0,c_Y=1,c_W=m,c_Z=1+m$, ordered increasingly when applying the
theorem.  Write the trace fields as $x,y,w,z,D$.  At $m=1$ the clocks
are $0,1,2$, the fields are $x,q,z,D$, and $q=a_U$; notation such as
$y+w$ means $q$ and asserts no separate traces for $Y,W$.
Thus \eqref{eq:trace-distribution-inequality} reads
\begin{equation}\label{eq:trace-z-D}
 (\partial_u-(1+m)\partial_i)z\ge\partial_iD.
\end{equation}
The distinction between grouped and separate depths is necessary:
$f(L_Y,L_W)=\max(L_Y,L_W)$ has nonpositive mixed derivatives, but its
separate first derivatives have different traces from the two sides
of $L_Y=L_W$.  On the grouped domain it is simply $f(L_U)=L_U$.

\subsection{Traces on a contact curve}

We need a second trace statement within the plane.  It uses bounded
variation in one direction at a time; it does not require the entire
gradient of each rate field to be a measure.

\begin{lemma}[Transverse traces on a Lipschitz curve]
\label{lem:curve-traces}
Let $b$ be Lipschitz on an interval $J$, and let
$\Gamma=\{(b(u),u):u\in J\}$ lie in an open plane region.
Suppose bounded fields $a_j,D$ in this region satisfy
$V_ja_j\in\mathcal M_{\rm loc}$ and
$\partial_iD\in\mathcal M_{\rm loc}$, where $\mathcal M_{\rm loc}$
denotes locally finite signed measures. Put~$p(u)=-b'(u)$ at points
where this derivative exists.  For almost every $u\in J$ the field
$D$ has two strong traces $D^-,D^+$ and every $a_j$ with
$c_j\ne p(u)$ has two strong traces $a_j^-,a_j^+$.  Specifically, at
$q_0=(b(u_0),u_0)$ the mean of
$|a_j-a_j^\pm(u_0)|$ tends to zero on the corresponding tangent
half-ball
\[
 B_\rho(q_0)\cap
 \{\ \pm(i-b(u_0)+p(u_0)(u-u_0))>0\ \},
 \qquad\rho\downarrow0,
\]
and likewise for $D$.  The same assertion holds with the actual
sides $\pm(i-b(u))>0$.  No assertion is made about the field whose
speed equals $p(u_0)$.

If, in addition, $V_na_n\ge\partial_iD$, then at almost every
transverse point $p\ne c_n$,
\begin{equation}\label{eq:trace-curve-jump}
 D^- -D^+\ge(c_n-p)(a_n^+-a_n^-).
\end{equation}
These conclusions apply to curves on the boundary of an aperture
region whenever the fields are defined in an ambient neighborhood;
one then retains only the side lying in that region.
\end{lemma}

\begin{proof}
We first justify the one-direction slicing fact being used.  If
$g\in L^\infty$ on a rectangle with coordinates $(s,r)$ and
$\partial_sg$ is a finite measure, then for almost every $r$ the
one-dimensional function $g_r(s)=g(s,r)$ has bounded variation, and
its derivative measures satisfy, on smaller rectangles,
\begin{equation}\label{eq:trace-slicing}
 \partial_sg=\int Dg_r\,dr,
 \qquad |\partial_sg|=\int |Dg_r|\,dr.
\end{equation}
Here integration of measures means integration of their values
against compactly supported continuous test functions.  To verify
this, convolve $g$ in the rectangle.  The integrals of
$|\partial_sg_\varrho|$ on a smaller rectangle are bounded by the
variation of $\partial_sg$ on a slightly larger one.  Choose a
subsequence with $g_\varrho(\cdot,r)\to g(\cdot,r)$ in $L^1(ds)$ for
almost every $r$.  The definition of one-dimensional variation as
the supremum over finite partitions, or over test functions bounded
by one, and Fatou's lemma give integrable variations of the slices.
Fubini and one-dimensional integration by parts then identify
$\partial_sg$ with $\int Dg_r\,dr$.  The variation equality follows
by taking the supremum over a countable dense family of scalar test
functions on rational subintervals, selecting these tests measurably
in $r$, and approximating the resulting bounded tests by continuous
ones.  Exhaustion gives \eqref{eq:trace-slicing} locally.

Suppose first the curve is a Lipschitz graph $s=h(r)$ in these
coordinates.  The one-dimensional representatives have limits
$T^\pm(r)=g_r(h(r)\pm)$ for almost every $r$.  Fix a compact graph
patch and define
\[
 q_\varrho(r)=|Dg_r|\bigl((h(r)-\varrho,h(r)+\varrho)
                                      \setminus\{h(r)\}\bigr).
\]
For small fixed $\varrho$, these functions are integrable by
\eqref{eq:trace-slicing}, and they decrease to zero almost everywhere
as $\varrho\downarrow0$.  Outside one null set we may assume that
$r_0$ is a Lebesgue point of $T^+,T^-$ and of every $q_{1/k}$ for
which the graph patch permits that collar.  If
$|r-r_0|<C\rho$ and $0<\pm(s-h(r))<C\rho$, then
\[
 |g(s,r)-T^\pm(r_0)|
 \le q_{C\rho}(r)+|T^\pm(r)-T^\pm(r_0)|
\]
for almost every $(s,r)$.  Integrating in such a graph-side cylinder
and dividing by $\rho^2$ bounds the result by
\[
 \frac{C}{\rho}\int_{|r-r_0|<C\rho}
       \bigl(q_{C\rho}(r)+|T^\pm(r)-T^\pm(r_0)|\bigr)\,dr.
\]
For each fixed $k$, eventually $q_{C\rho}\le q_{1/k}$.
Lebesgue differentiation first, then $k\to\infty$, shows that this
bound tends to zero.  At a point where $h$ is differentiable, the
area between its graph and its tangent in a radius-$\rho$ disk is
$o(\rho^2)$.  Boundedness of $g$ therefore gives the same strong mean
traces on tangent half-disks.  Invertible linear changes of coordinates
preserve this conclusion for half-balls, because the corresponding
ellipses contain and are contained in comparable balls.

For $a_j$ use the coordinates
\[
 s=u,\qquad r=i+c_ju,
\]
so that $V_j=\partial_s$ and the curve has projection
$r=g_j(u)=b(u)+c_ju$.  On the set where $g_j'=c_j-p\ne0$ this
projection can be covered, apart from a null set, by countably many
sets on each of which it is bi-Lipschitz.  Here is a direct reason.
At a differentiability point with $g_j'\ne0$, choose integers $k,N$
so that $|g_j'|>2/k$ and
\[
 |g_j(v)-g_j(u)-g_j'(u)(v-u)|\le |v-u|/k
 \quad\text{if }|v-u|<1/N.
\]
Partition these points according to $k,N$, the sign of the derivative,
and intervals of length less than $1/N$.  Any two points of one piece
then have image distance at least $|v-u|/k$; the upper Lipschitz bound
is inherited from $g_j$.  On each image piece the inverse extends
to a Lipschitz function $h(r)$ on an interval.  The graph argument
above applies to this extension.  At almost every density point of
the original piece its tangent is the original tangent.  The
bi-Lipschitz parameter change takes exceptional null sets back to
null sets.  This proves the traces for $a_j$ on its transverse set.
For $D$, simply take $s=i,r=u$, so the curve is already the graph
$s=b(r)$ and is everywhere transverse.  Taking the union of the
finitely many exceptional sets proves the simultaneous assertion.
The area comparison also replaces tangent sides by actual sides.

Finally, \eqref{eq:trace-slicing} identifies the part of the derivative
measure on $s=h(r)$ as $(T^+(r)-T^-(r))\,dr$, since a
one-dimensional BV derivative has an atom equal to its jump.
On each transverse inverse patch, substitute
$dr=|c_j-p(u)|\,du$.  The positive $s$ side is the positive $i-b(u)$
side precisely when $p-c_j>0$, so reversing the sides when necessary
gives the signed factor $p-c_j$.  Partition the countable patches
into disjoint measurable pieces before summing.  There is no omitted
measure on a parameter-null subset of the curve: its projection under
the Lipschitz map $g_j$ is null, and \eqref{eq:trace-slicing} gives zero
variation over such a projection.  In the original $(i,u)$ coordinates
we have therefore proved
\begin{equation}\label{eq:trace-graph-measures}
 (V_ja_j)|_\Gamma
   =(p-c_j)(a_j^+-a_j^-)\,du,
 \qquad
 (\partial_iD)|_\Gamma=(D^+-D^-)\,du
\end{equation}
on the transverse part of the parameter set.  To check the first
sign and factor directly, set $\varphi(i,u)=i-b(u)$.  Then
$V_j\varphi=p-c_j$, and differentiating a function with jump
$a_j^+-a_j^-$ across $\varphi=0$ gives exactly that factor times
$\delta(\varphi)\,di\,du=(du)|_\Gamma$.
This agrees with the change of variable in the sliced jump measures,
so it applies to the bounded fields just considered.  Restrict the
nonnegative measure $V_na_n-\partial_iD$ to the transverse part of
$\Gamma$ and use
\eqref{eq:trace-graph-measures}; its nonnegative density is
\[
 (p-c_n)(a_n^+-a_n^-)-(D^+-D^-).
\]
This is precisely \eqref{eq:trace-curve-jump}.  An aperture boundary
requires no further trace argument: all statements were proved in
the ambient plane neighborhood before a side is selected.
\end{proof}
